\chapter{Real-парный радиационный ток и родитель внутреннего охлаждения}
\label{ch:v6-spectral-transition-real-pair-radiative-cooling-parent-gate}

\section{Решающий вопрос}

Переприоритизация оставила один минимальный мост: использовать точный
исходящий канал compacton не для выбора частицы, а как возможный перенос
энергии и энтропии проекторной кристаллизации. Первый стоп-критерий
формулируется до любой термодинамической интерпретации: переживает ли
радиационный ток полную Real/KO6-пару или сокращается вместе с
ориентационной фазой.

Если положительный ток сохраняется, необходимо проверить более сильное
условие: существуют ли уже в проекте общий носитель, взаимодействие и
нормировка, превращающие норму исходящего пакета в энергетический баланс
параметра порядка $R$.

\section{Радиационные пакеты двух сопряжённых ветвей}

Обозначим compacton-ветви через $\Psi_+$ и $\Psi_-$,
\begin{equation}
 F(\Psi_\pm)=\pm i\Psi_\pm,
 \qquad
 \Psi_-=J\Psi_+.
\end{equation}
Пусть $\mathcal M_\pm$ --- вещественная производная четырёхшаговой карты
около соответствующей ветви. Для сопряжённых характерных возмущений
определим исходящие части
\begin{align}
 r_+&=P_{\rm out}\bigl(\mathcal M_+\Psi_- -\Psi_-\bigr),\\
 r_-&=P_{\rm out}\bigl(\mathcal M_-\Psi_+ -\Psi_+\bigr).
 \label{eq:v6-real-pair-radiation-packets}
\end{align}

Машинный расчёт на $256$ узлах даёт
\begin{equation}
 \|r_+-\overline{r_-}\|=2.63\cdot10^{-10},
 \qquad
 \|r_+\|^2=\|r_-\|^2
 =39.4784176072\ldots.
\end{equation}
Следовательно, линейная ориентированная величина сокращается при
сопряжении, но положительный физический полуслед равен
\begin{equation}
 J_{\rm rad}^{\rm Real}
 =\frac12\left(\|r_+\|^2+\|r_-\|^2\right)
 =4\pi^2
 \label{eq:v6-real-pair-positive-radiation-current}
\end{equation}
с остатком $2.87\cdot10^{-9}$.

\begin{proposition}[Положительный ток не сокращается]
Real-сопряжение уничтожает ориентационно нечётную фазовую амплитуду, но
не уничтожает квадратичный радиационный ток. После нормированного
физического полуследа коэффициент остаётся равным $4\pi^2$: он не
удваивается и не обращается в ноль.
\end{proposition}

Этот результат согласуется с общей структурой Real-пары: суперслед или
фаза могут сокращаться, тогда как положительная норма сохраняется.

\section{Условное энергетическое сопоставление}

Проекторный переход к состоянию сосуществования высвобождает в принятой
безразмерной нормировке
\begin{equation}
 \Delta E_R=0.4798400112\ldots.
\end{equation}
Если временно, только для диагностики, отождествить единицу касательной
нормы блуждания с единицей проекторной энергии, то за $N$ циклов баланс
требовал бы
\begin{equation}
 N\,4\pi^2|\delta|^2=\Delta E_R.
 \label{eq:v6-real-pair-conditional-energy-balance}
\end{equation}
При $N=1,2,4,8,16,32$ условные амплитуды равны соответственно
\begin{equation}
 0.11025, 0.07796, 0.05512, 0.03898, 0.02756, 0.01949.
\end{equation}
Эта таблица не является предсказанием. В общем случае требуется
неизвестная конверсионная нормировка $\chi$,
\begin{equation}
 \chi N\,4\pi^2|\delta|^2=\Delta E_R.
 \label{eq:v6-real-pair-free-conversion}
\end{equation}
Одно равенство не фиксирует $\chi$, $N$ и $\delta$. Кроме того, физическое
время равно $4N\Delta t$, а $\Delta t$ ранее осталось свободным.

Иными словами, совпадение безразмерных чисел в двух разных действиях не
создаёт общего энергетического баланса. Коэффициент $\chi=1$ допустим
только как следствие единой родительской нормы, а не как соглашение.

\section{Почему ортогональные пакеты ещё не являются стоком энтропии}

Глобальное нелинейное блуждание сохраняет чистоту и норму. Появление
пространственно ортогональных пакетов даёт память о циклах, но само по себе
не увеличивает энтропию глобального состояния. Термодинамический экспорт
возникает лишь для редуцированного состояния после физически заданного
разложения
\begin{equation}
 \mathcal H_{\rm parent}
 =\mathcal H_R\otimes\mathcal H_{\rm out}otimes\mathcal H_{\rm aux}
\end{equation}
и частичного следа по исходящим модам.

Такая редукция не является произвольной математической операцией: она
должна следовать из локального взаимодействия, причинного удаления пакетов
или термодинамического предела. Автономные часы и холодильники требуют
явного энергетического носителя и испытывают обратную реакцию
\cite{real-pair-cooling-erker2017,real-pair-cooling-linden2010}.
Ненулевая ширина в непрерывный канал задаёт распад только после
идентификации дискретной моды, континуума и их общего гамильтониана
\cite{real-pair-cooling-fano1961}.

В текущей конструкции compacton-блуждание определено на пространственно-
хиральном секторе, тогда как $R$ живёт в семейном проекторном секторе.
Общий тензорный носитель, взаимодействие и частичный след между ними не
предъявлены. Поэтому требуемый экспорт
\begin{equation}
 \Delta S_R=0.7402345698\ldots
\end{equation}
не следует из одной ортогональности пакетов.

\section{Родительский ledger}

Из семи обязательных пунктов проходит только первый:
\begin{center}
\begin{tabular}{p{0.67\textwidth}p{0.16\textwidth}}
\toprule
Условие & Статус \\
\midrule
Real-чётный поток переживает полуслед & пройдено \\
Общий носитель $R$ и compacton-блуждания & не выведен \\
Общая нормировка радиационной и проекторной энергии & не выведена \\
Родительское взаимодействие $R$ с исходящими модами & не выведено \\
Амплитуда $\delta$ следует из состояния & нет \\
Монотонный закон $\beta(\tau)$ & нет \\
Физический частичный след, задающий сток энтропии & не выведен \\
\bottomrule
\end{tabular}
\end{center}

Каноническая аффинная связь $\rho V$ не закрывает пробел: она действует
скалярно на семейном триплете и сохраняет $R=I_3/3$. Чтобы радиация
изменила семейный спектр, требуется общий оператор, а не только совпадение
размерностей или характеров.

\begin{nogocriterion}[Норма пакета не равна теплу]
Нельзя положить $\chi=1$ в
\eqref{eq:v6-real-pair-free-conversion}, назвать исходящие моды
резервуаром и затем вывести $\beta(\tau)$. Пока общая мера и
взаимодействие отсутствуют, это три независимых назначения: единицы
энергии, амплитуды возбуждения и правила крупнозернения.
\end{nogocriterion}

\begin{theorem}[Частичный успех и провал родительского моста]
Положительный радиационный ток полной Real-пары ненулевой и имеет точный
коэффициент $4\pi^2$. Тем самым первое предполагаемое сокращение
опровергнуто. Однако существующая версия не содержит общего носителя,
энергетической нормировки, взаимодействия и редуцированного следа,
связывающих этот ток с проекторным параметром порядка. Амплитуда
$\delta$, шаг времени и закон $\beta(\tau)$ остаются свободными.
Следовательно, радиационный канал не является выведенным внутренним
охлаждением.
\end{theorem}

\section{Следующий гейт}

Следующий и последний минимальный тест этого маршрута
\path{version6_spectral_transition_radiative_cooling_common_carrier_attribution_gate}
должен проверить, существует ли уже в полном носителе
$M_{300}(\mathbb C)$ канонический общий угол или суперсвязностная стрелка,
которая одновременно содержит семейный $R$, хиральный compacton-сектор и
исходящие пространственные моды с одной следовой нормировкой.

Нулевая гипотеза: существующие связи факторизуются как тождество на
семейном секторе и не меняют $R=I_3/3$. Стоп-критерий: если общий
intertwiner отсутствует или определён только с новым коэффициентом,
радиационно-охлаждающий маршрут окончательно закрывается.

\begin{protocol}
\begin{enumerate}
 \item сопряжённость исходящих пакетов: \textbf{подтверждена};
 \item ориентационно нечётная амплитуда Real-пары: \textbf{сокращается};
 \item положительный полуслед потока: \textbf{$4\pi^2$};
 \item остаток от $4\pi^2$: \textbf{$2.87\cdot10^{-9}$};
 \item общий носитель с проекторным $R$: \textbf{не предъявлен};
 \item конверсионная нормировка $\chi$: \textbf{не выведена};
 \item амплитуда $\delta$: \textbf{не выведена};
 \item физический шаг времени: \textbf{не выведен};
 \item экспорт энтропии одной унитарной радиацией: \textbf{нет};
 \item монотонный $\beta(\tau)$: \textbf{не выведен};
 \item родительский тест: \textbf{1 из 7 пунктов};
 \item Real-парное охлаждение: \textbf{не доказано};
 \item сохранённая зацепка: \textbf{ненулевой положительный ток}.
\end{enumerate}
\end{protocol}

Машинный сертификат сохранён в
\path{s2t_v6_spectral_transition_real_pair_radiative_cooling_parent_gate_results.json}.

\begin{thebibliography}{9}
\bibitem{real-pair-cooling-erker2017}
P.~Erker, M.~T.~Mitchison, R.~Silva, M.~P.~Woods, N.~Brunner, and
M.~Huber, \emph{Autonomous Quantum Clocks: Does Thermodynamics Limit Our
Ability to Measure Time?}, Physical Review X 7 (2017), 031022,
\path{arXiv:1609.06704}.

\bibitem{real-pair-cooling-linden2010}
N.~Linden, S.~Popescu, and P.~Skrzypczyk,
\emph{How Small Can Thermal Machines Be? The Smallest Possible
Refrigerator}, Physical Review Letters 105 (2010), 130401,
\path{doi:10.1103/PhysRevLett.105.130401}.

\bibitem{real-pair-cooling-fano1961}
U.~Fano,
\emph{Effects of Configuration Interaction on Intensities and Phase Shifts},
Physical Review 124 (1961), 1866--1878,
\path{doi:10.1103/PhysRev.124.1866}.
\end{thebibliography}